\section{Proper morphisms}
\label{section-proper}

\noindent
The notion of a proper morphism plays an important role in algebraic
geometry. An important example of a proper morphism will be the
structure morphism $\mathbf{P}^n_S \to S$ of projective $n$-space,
and this is in fact the motivating example leading to the definition.

\begin{definition}
\label{definition-proper}
Let $f : X \to S$ be a morphism of schemes.
We say $f$ is {\it proper} if $f$ is separated, finite type, and
universally closed.
\end{definition}

\noindent
The morphism from the affine line with zero doubled to the affine line
is of finite type and universally closed, so the separation condition is
necessary in the definition above.
In the rest of this section we prove some of the basic properties
of proper morphisms and of universally closed morphisms.

\begin{lemma}
\label{lemma-universally-closed-local-on-the-base}
Let $f : X \to S$ be a morphism of schemes.
The following are equivalent:
\begin{enumerate}
\item The morphism $f$ is universally closed.
\item There exists an open covering $S = \bigcup V_j$ such
that $f^{-1}(V_j) \to V_j$ is universally closed for all indices $j$.
\end{enumerate}
\end{lemma}

\begin{proof}
This is clear from the definition.
\end{proof}

\begin{lemma}
\label{lemma-proper-local-on-the-base}
Let $f : X \to S$ be a morphism of schemes.
The following are equivalent:
\begin{enumerate}
\item The morphism $f$ is proper.
\item There exists an open covering $S = \bigcup V_j$ such
that $f^{-1}(V_j) \to V_j$ is proper for all indices $j$.
\end{enumerate}
\end{lemma}

\begin{proof}
Omitted.
\end{proof}

\begin{lemma}
\label{lemma-composition-proper}
The composition of proper morphisms is proper.
The same is true for universally closed morphisms.
\end{lemma}

\begin{proof}
A composition of closed morphisms is closed.
If $X \to Y \to Z$ are universally closed morphisms
and $Z' \to Z$ is any morphism, then we see that
$Z' \times_Z X = (Z' \times_Z Y) \times_Y X  \to Z' \times_Z Y$
is closed and $Z' \times_Z Y \to Z'$ is closed.
Hence the result for universally closed morphisms.
We have seen that ``separated'' and ``finite type''
are preserved under compositions
(Schemes, Lemma \ref{schemes-lemma-separated-permanence} and
Lemma \ref{lemma-composition-finite-type}). Hence the
result for proper morphisms.
\end{proof}

\begin{lemma}
\label{lemma-base-change-proper}
The base change of a proper morphism is proper.
The same is true for universally closed morphisms.
\end{lemma}

\begin{proof}
This is true by definition for universally closed morphisms.
It is true for separated morphisms
(Schemes, Lemma \ref{schemes-lemma-separated-permanence}).
It is true for morphisms of finite type
(Lemma \ref{lemma-base-change-finite-type}).
Hence it is true for proper morphisms.
\end{proof}

\begin{lemma}
\label{lemma-closed-immersion-proper}
A closed immersion is proper, hence a fortiori universally closed.
\end{lemma}

\begin{proof}
The base change of a closed immersion is a closed immersion
(Schemes, Lemma \ref{schemes-lemma-base-change-immersion}).
Hence it is universally closed.
A closed immersion is separated
(Schemes, Lemma \ref{schemes-lemma-immersions-monomorphisms}).
A closed immersion is of finite type
(Lemma \ref{lemma-immersion-locally-finite-type}).
Hence a closed immersion is proper.
\end{proof}

\begin{lemma}
\label{lemma-image-proper-scheme-closed}
Suppose given a commutative diagram of schemes
$$
\xymatrix{
X \ar[rr] \ar[rd] & &
Y \ar[ld] \\
& S &
}
$$
with $Y$ separated over $S$.
\begin{enumerate}
\item If $X \to S$ is universally closed, then the morphism
$X \to Y$ is universally closed.
\item If $X$ is proper over $S$, then the morphism $X \to Y$ is proper.
\end{enumerate}
In particular, in both cases the image of $X$ in $Y$ is closed.
\end{lemma}

\begin{proof}
Assume that $X \to S$ is universally closed (resp.\ proper).
We factor the morphism as $X \to X \times_S Y \to Y$.
The first morphism is a closed immersion, see
Schemes, Lemma \ref{schemes-lemma-semi-diagonal}.
Hence the first morphism is proper (Lemma \ref{lemma-closed-immersion-proper}).
The projection $X \times_S Y \to Y$ is the base change
of a universally closed (resp.\ proper) morphism and hence
universally closed (resp.\ proper), see Lemma \ref{lemma-base-change-proper}.
Thus $X \to Y$ is universally closed (resp.\ proper) as the composition
of universally closed (resp.\ proper) morphisms
(Lemma \ref{lemma-composition-proper}).
\end{proof}

\noindent
The proof of the following lemma is due to Bjorn Poonen, see
\href{https://mathoverflow.net/questions/23337/is-a-universally-closed-morphism-of-schemes-quasi-compact/23528#23528}{this location}.

\begin{lemma}
\label{lemma-universally-closed-quasi-compact}
\begin{reference}
Due to Bjorn Poonen.
\end{reference}
A universally closed morphism of schemes is quasi-compact.
\end{lemma}

\begin{proof}
Let $f : X \to S$ be a morphism. Assume that $f$ is not quasi-compact.
Our goal is to show that $f$ is not universally closed. By
Schemes, Lemma \ref{schemes-lemma-quasi-compact-affine}
there exists an affine open $V \subset S$ such that $f^{-1}(V)$ is
not quasi-compact. To achieve our goal it suffices to show that
$f^{-1}(V) \to V$ is not universally closed, hence we may assume that
$S = \Spec(A)$ for some ring $A$.

\medskip\noindent
Write $X = \bigcup_{i \in I} X_i$ where the $X_i$ are affine open subschemes
of $X$.  Let $T = \Spec(A[y_i ; i \in I])$.
Let $T_i = D(y_i) \subset T$. Let $Z$ be the closed set
$(X \times_S T) - \bigcup_{i \in I} (X_i \times_S T_i)$.  It suffices to
prove that the image $f_T(Z)$ of $Z$ under $f_T : X \times_S T \to T$
is not closed.

\medskip\noindent
There exists a point $s \in S$ such that there is no
neighborhood $U$ of $s$ in $S$ such that $X_U$ is quasi-compact.
Otherwise we could cover $S$ with finitely many such $U$ and
Schemes, Lemma \ref{schemes-lemma-quasi-compact-affine}
would imply $f$ quasi-compact. Fix such an $s \in S$.

\medskip\noindent
First we check that $f_T(Z_s) \ne T_s$.  Let $t \in T$ be the point
lying over $s$ with $\kappa(t) = \kappa(s)$ such that $y_i = 1$ in
$\kappa(t)$ for all $i$.
Then $t \in T_i$ for all $i$, and the fiber of $Z_s \to T_s$ above
$t$ is isomorphic to $(X - \bigcup_{i \in I} X_i)_s$, which is empty.
Thus $t \in T_s - f_T(Z_s)$.

\medskip\noindent
Assume $f_T(Z)$ is closed in $T$. Then there exists an element
$g \in A[y_i; i \in I]$ with $f_T(Z) \subset V(g)$ but $t \not \in V(g)$.
Hence the image of $g$ in $\kappa(t)$ is nonzero. In particular some
coefficient of $g$ has nonzero image in $\kappa(s)$. Hence this coefficient is
invertible on some neighborhood $U$ of $s$. Let $J$ be the finite set of
$j \in I$ such that $y_j$ appears in $g$. Since $X_U$ is not quasi-compact,
we may choose a point $x \in X - \bigcup_{j \in J} X_j$ lying above some
$u \in U$. Since $g$ has a coefficient that is invertible on $U$, we can
find a point $t' \in T$ lying above $u$ such that $t' \not \in V(g)$ and
$t' \in V(y_i)$ for all $i \notin J$. This is true because
$V(y_i; i \in I, i \not\in J) = \Spec(A[y_j; j\in J])$
and the set of points of this scheme lying over $u$ is bijective
with $\Spec(\kappa(u)[y_j; j \in J])$. In other words $t' \notin T_i$
for each $i \notin J$. By
Schemes, Lemma \ref{schemes-lemma-points-fibre-product}
we can find a point $z$ of $X \times_S T$ mapping to $x \in X$ and to
$t' \in T$. Since $x \not \in X_j$ for $j \in J$ and $t' \not \in T_i$
for $i \in I \setminus J$ we see that $z \in Z$. On the other hand
$f_T(z) = t' \not \in V(g)$ which contradicts $f_T(Z) \subset V(g)$.
Thus the assumption ``$f_T(Z)$ closed'' is wrong and we conclude indeed
that $f_T$ is not closed, as desired.
\end{proof}

\noindent
The following lemma says that the image of a proper scheme (in a separated
scheme of finite type over the base) is proper.

\begin{lemma}
\label{lemma-image-proper-is-proper}
Let $S$ be a scheme. Let $f : X \to Y$ be a morphism of schemes over $S$.
If $X$ is universally closed over $S$ and $f$ is surjective then
$Y$ is universally closed over $S$. In particular, if also $Y$ is
separated and locally of finite type over $S$, then $Y$ is proper over $S$.
\end{lemma}

\begin{proof}
Assume $X$ is universally closed and $f$ surjective.
Denote $p : X \to S$, $q : Y \to S$ the structure morphisms.
Let $S' \to S$ be a morphism of schemes. The base change
$f' : X_{S'} \to Y_{S'}$ is surjective
(Lemma \ref{lemma-base-change-surjective}), and the base
change $p' : X_{S'} \to S'$ is closed.
If $T \subset Y_{S'}$ is closed, then $(f')^{-1}(T) \subset X_{S'}$
is closed, hence $p'((f')^{-1}(T)) = q'(T)$ is closed.
So $q'$ is closed. This proves the first statement.
Thus $Y \to S$ is quasi-compact by
Lemma \ref{lemma-universally-closed-quasi-compact}
and hence $Y \to S$ is proper by definition
if in addition $Y \to S$ is locally of finite type and separated.
\end{proof}

\begin{lemma}
\label{lemma-scheme-theoretic-image-is-proper}
Suppose given a commutative diagram of schemes
$$
\xymatrix{
X \ar[rr]_h \ar[rd]_f & & Y \ar[ld]^g \\
& S
}
$$
Assume
\begin{enumerate}
\item $X \to S$ is a universally closed (for example proper) morphism, and
\item $Y \to S$ is separated and locally of finite type.
\end{enumerate}
Then the scheme theoretic image $Z \subset Y$ of $h$
is proper over $S$ and $X \to Z$ is surjective.
\end{lemma}

\begin{proof}
The scheme theoretic image of $h$ is constructed in Section
\ref{section-scheme-theoretic-image}. Since $f$ is quasi-compact
(Lemma \ref{lemma-universally-closed-quasi-compact})
we find that $h$ is quasi-compact
(Schemes, Lemma \ref{schemes-lemma-quasi-compact-permanence}).
Hence $h(X) \subset Z$
is dense (Lemma \ref{lemma-quasi-compact-scheme-theoretic-image}).
On the other hand $h(X)$ is closed in $Y$
(Lemma \ref{lemma-image-proper-scheme-closed})
hence $X \to Z$ is surjective.
Thus $Z \to S$ is a proper (Lemma \ref{lemma-image-proper-is-proper}).
\end{proof}

\noindent
The target of a separated scheme under a surjective
universally closed morphism is separated.

\begin{lemma}
\label{lemma-image-universally-closed-separated}
Let $S$ be a scheme. Let $f : X \to Y$ be a surjective universally closed
morphism of schemes over $S$.
\begin{enumerate}
\item If $X$ is quasi-separated, then $Y$ is quasi-separated.
\item If $X$ is separated, then $Y$ is separated.
\item If $X$ is quasi-separated over $S$, then $Y$ is quasi-separated over $S$.
\item If $X$ is separated over $S$, then $Y$ is separated over $S$.
\end{enumerate}
\end{lemma}

\begin{proof}
Parts (1) and (2) are a consequence of (3) and (4) for
$S = \Spec(\mathbf{Z})$ (see
Schemes, Definition \ref{schemes-definition-separated}).
Consider the commutative diagram
$$
\xymatrix{
X \ar[d] \ar[rr]_{\Delta_{X/S}} & & X \times_S X \ar[d] \\
Y \ar[rr]^{\Delta_{Y/S}} & & Y \times_S Y
}
$$
The left vertical arrow is surjective (i.e., universally surjective).
The right vertical arrow is universally closed as a composition
of the universally closed morphisms
$X \times_S X \to X \times_S Y \to Y \times_S Y$. Hence it is also
quasi-compact, see
Lemma \ref{lemma-universally-closed-quasi-compact}.

\medskip\noindent
Assume $X$ is quasi-separated over $S$, i.e.,  $\Delta_{X/S}$ is
quasi-compact. If $V \subset Y \times_S Y$ is a quasi-compact
open, then $V \times_{Y \times_S Y} X \to \Delta_{Y/S}^{-1}(V)$
is surjective and $V \times_{Y \times_S Y} X$ is quasi-compact by our remarks
above. We conclude that $\Delta_{Y/S}$ is quasi-compact, i.e., $Y$
is quasi-separated over $S$.

\medskip\noindent
Assume $X$ is separated over $S$, i.e., $\Delta_{X/S}$ is a closed
immersion. Then $X \to Y \times_S Y$ is closed as a
composition of closed morphisms. Since $X \to Y$ is
surjective, it follows that $\Delta_{Y/S}(Y)$ is closed in $Y \times_S Y$.
Hence $Y$ is separated over $S$ by the discussion following
Schemes, Definition \ref{schemes-definition-separated}.
\end{proof}








